% © 2026 Ing. David Jaroš — CC BY-NC-ND 4.0
%
% This work is licensed under a Creative Commons Attribution-NonCommercial-NoDerivatives
% 4.0 International License (CC BY-NC-ND 4.0).
%
% See LICENSE.md for full license text.

\documentclass[12pt,a4paper]{article}
\usepackage[a4paper,margin=2.5cm]{geometry}
\usepackage{amsmath,amssymb,tcolorbox}
\tcbuselibrary{skins,breakable}
\newtcolorbox{statusbox}[1][]{colback=gray!8!white,colframe=black!60,arc=3pt,boxrule=1pt,title=Status Summary,#1}

\title{Higher Harmonics Above the Three Fermion Generations}
\author{Ing.\ David Jaro\v{s}}
\date{2026-05-12}

\begin{document}
\maketitle

\section{Problem Statement}

UBT uses winding structure on compact phase circles to explain why the observed
fermions come in three generations.  The immediate consistency question is that
the winding spectrum on \(S^1\) is infinite:
\[
n\in\mathbb{Z}.
\]
If the first three light sectors are associated with the observed generations,
what should be done with the modes \(n\ge 4\)?

The issue becomes sharper in the SU(2)-twist route, where the charged
off-diagonal fields obey anti-periodic boundary conditions and therefore carry
the standard half-integer KK spectrum
\[
m_n=\frac{n+1/2}{R_\psi},\qquad n=0,1,2,\dots
\]
instead of the periodic tower \(m_n=n/R_\psi\).  In that notation the first
three candidate generation levels are the three lowest states
\[
\frac{1}{2R_\psi},\quad \frac{3}{2R_\psi},\quad \frac{5}{2R_\psi},
\]
which may be labeled generation \(1,2,3\) even though the underlying KK index is
\(n=0,1,2\).

\section{Candidate A: Massive KK Excitations}

This is the most natural interpretation of the higher harmonics.

For a periodic compact direction the KK formula is
\[
m_n=\frac{|n|}{R_\psi}.
\]
For the charged SU(2)-twisted sector used in
\texttt{research\_tracks/quantum\_ubt/su2\_twist\_neff12.tex}, the correct
formula is instead
\[
m_n^{\mathrm{twist}}=\frac{n+1/2}{R_\psi},
\]
because the fields \(b,c\) are anti-periodic while the neutral fields \(a,d\)
remain periodic.  This part is standard KK/Scherk--Schwarz theory once the
twisted boundary condition is assumed.

Hence the next charged levels after the first three are
\[
\frac{7}{2R_\psi},\ \frac{9}{2R_\psi},\ \frac{11}{2R_\psi},\dots
\]
and for \(R_\psi^{-1}\) near a UV scale (GUT or Planck-like), these states are
far beyond present experimental reach.  For the Planck estimate mentioned in the
problem statement, even the first omitted level is of order
\[
\frac{7}{2}M_{\mathrm{Pl}}\sim 4\times10^{19}\,\mathrm{GeV},
\]
so the higher harmonics simply decouple from low-energy phenomenology.

What this candidate \emph{does} explain is why \(n\ge 4\) modes are not seen.
What it does \emph{not} yet prove is why exactly the first three levels should be
identified with the observed light generations after all symmetry breaking and
mass generation effects are included.

\section{Candidate B: Algebraic Selection Rule}

There is currently no convincing algebraic rule in \(\mathbb{C}\otimes\mathbb{H}\)
that truncates the winding tower at three.

The algebra does contain a distinguished ``three'':
\[
\dim_{\mathbb{R}}(\mathrm{Im}\,\mathbb{H})=3.
\]
That supports the statement that there are three independent phase directions,
or three structurally distinguished channels.  But it does \emph{not} imply that
each compact circle carries only three Fourier modes.  Once a compact circle
\(S^1_\psi\) is present, the spectrum of winding numbers is automatically
infinite.

Therefore a purely algebraic cutoff \(n\ge 4\mapsto 0\) is not available from
the biquaternion algebra alone.  Candidate~B is therefore a \textbf{[NO-GO]} as
a standalone explanation.

\section{Candidate C: Composite States}

The composite-state interpretation is logically possible but currently
unsupported.

Mode numbers on a compact circle add under momentum conservation, so one can
write schematic relations such as \(4=1+3\).  But that by itself does not make
the \(n=4\) harmonic a bound state of the \(n=1\) and \(n=3\) harmonics.  To
justify a genuinely composite interpretation one would need:
\begin{itemize}
  \item an interaction term in \(S[\Theta]\) that binds lower modes into higher
  resonances,
  \item a selection rule separating fundamental and composite towers, and
  \item a mass/width analysis showing resonance behaviour rather than a single KK
  eigenstate.
\end{itemize}
None of these ingredients has been derived in the active EW documents.  The
composite interpretation remains \textbf{[OPEN]} and presently weaker than the
KK explanation.

\section{Verdict}

The best current interpretation is:
\[
\boxed{\text{higher harmonics }n\ge 4\text{ are massive KK excitations.}}
\]

Candidate~A fits the active compact-circle formalism, matches the standard
twisted KK formula, and explains why no extra generations appear at accessible
energies.  Candidate~B fails because the algebraic ``three'' counts phase
directions, not Fourier harmonics.  Candidate~C is not excluded in principle,
but no derivation supports it.

The remaining open point is not the existence of higher modes; it is the
low-energy selection problem.  UBT still needs a first-principles derivation of
why precisely the first three winding sectors become the observed light fermion
families while the rest stay at the compactification scale.

\begin{statusbox}
Vyšší harmonické \(n\ge 4\): \textbf{MASSIVE KK}\\
Klasifikace: \textbf{[MC]} --- KK interpretace je nejlépe podložená aktivními soubory, ale úplná first-principles selekce přesně tří lehkých generací ještě není uzavřena.\\
Dopad na počet generací: \textbf{Nízká energie stále vidí tři lehké generace, pokud \(R_\psi^{-1}\) leží vysoko; vyšší módy rozšiřují spektrum jen jako těžké excitace, ne jako nové dostupné SM rodiny.}
\end{statusbox}

\end{document}
